\fichatitulo{28 · ARQUITETURA}{Preprint v1 · 2024}{Modular RAG}
{\small Gao et al. \textit{Modular RAG: Transforming RAG Systems into LEGO-like Reconfigurable Frameworks}. Preprint v1 · 2024. \href{https://arxiv.org/pdf/2407.21059v1}{Fonte consultada}.\par}

\campo{Problema e objetivo}
Como representar pipelines RAG além de uma sequência linear fixa?

\campo{Principal contribuição · síntese interpretativa}
Decompõe sistemas em módulos e operadores, com mecanismos explícitos de roteamento, agendamento e fusão.

\campo{Método / natureza da evidência}
Estrutura conceitual para organizar padrões de execução e combinações de componentes.

\campo{Resultado ou argumento principal}
Distingue formas lineares, condicionais, ramificadas e iterativas de recuperação e geração.

\campo{Excertos-chave · seleção literal}
\begin{itemize}
\excerto{1}{routing, scheduling, and fusion mechanisms}{Resumo.}
\excerto{2}{linear, conditional, branching, and looping}{Resumo.}
\end{itemize}

\campo{Limitações e análise crítica}
Análise da ficha: a organização conceitual não comprova ganho automático ao adicionar módulos. Mais etapas podem introduzir custo e dificultar atribuir resultados a um componente.

\campo{Aproveitamento na dissertação · proposta}
Metodologia: descrever claramente o pipeline e planejar ablações para cada alteração relevante.

\registro{Fonte consultada nesta preparação: Resumo e introdução. Chave: \texttt{gao2024modularRag}.}
